% Generated draft: requires mathematical and source review.
\documentclass[11pt]{article}
\usepackage[margin=25mm]{geometry}
\usepackage{fontspec}
\setmainfont{Latin Modern Roman}
\usepackage{amsmath,amssymb,mathtools,unicode-math}
\setmathfont{Latin Modern Math}
\usepackage{enumitem,needspace,xurl,hyperref}
\setlength{\parindent}{0pt}
\setlength{\parskip}{5pt}
\setlength{\emergencystretch}{3em}
\newcommand{\N}{\mathbb{N}}
\newcommand{\Z}{\mathbb{Z}}
\newcommand{\Q}{\mathbb{Q}}
\newcommand{\R}{\mathbb{R}}
\newcommand{\C}{\mathbb{C}}
\newcommand{\F}{\mathbb{F}}
\newcommand{\A}{\mathbb{A}}
\newcommand{\PP}{\mathbb{P}}
\newcommand{\E}{\mathbb{E}}
\newcommand{\eps}{\varepsilon}
\newcommand{\dd}{\,\mathrm{d}}
\begin{document}
\section*{UnsolvedMath ID: 20001950}
\textbf{Source problem code:} catalogue entry
\subsection{Definitions and mathematical statement}
This record owns historical Comments 1.3--1.5 following Conjecture 1.1, rather than the conjecture itself. Its initial fragment belongs to Comment 1.2, and its final heading starts the next section. The missing hypotheses of Conjecture 1.1 and the missing Kirwan manuscript prevent a stand-alone universal surjectivity conjecture from being recovered here.

Comment 1.3 says that for an abelian group, workshop participants tentatively thought Kirwan's original symplectic surjectivity argument could be repeated with the hyperkähler moment map under the preceding conjecture's hypotheses. They explicitly did not check the details. Comment 1.4 asserts that Kirwan's third hypothesis, confinement of the relevant gradient trajectories to compact sets, holds for linear \(G\)-actions on \(T^*\C^n\), even for nonabelian \(G\), by an argument analogous to [Sja98] for \(\C^n\). Its stated quiver-variety surjectivity consequence remains conditional on the conjecture.

Comment 1.5 concerns the singular complex moment-map level, conventionally written \(\mu_{\C}^{-1}(0)\); the extraction of that notation is damaged. It questions the use of the Carrell--Goresky plus construction [CG83] in Section 5 of Kirwan's manuscript because its cited results are homological and not stated equivariantly, and because the discussed flows appeared to supply an \(\R^2\)-action rather than a \(\C^*\)-action. S. Wu proposed that a filtration could support both homology and cohomology spectral sequences; S. Wu and J. Weitsman discussed the action issue without establishing applicability. These are historical doubts and tentative observations.

The supplied aid reports the following conditional bridge package, not a verification of that singular-level application. For an inclusion \(i:Y\hookrightarrow X\) to which integral singular homology and its natural universal coefficient sequences apply, assume
\[
 i_*:H_q(Y;\Z)\longrightarrow H_q(X;\Z)
\]
is a split injection for every \(q\). The reported conclusion is that \(i^*:H^q(X;\Z)\to H^q(Y;\Z)\) is surjective for every \(q\). A mere comparison of homology ranks is not the stated hypothesis.

For the equivariant extension assume a compact connected group \(G\) acts, the inclusion is equivariant, and put
\[
 R=H^*(BG;\Z),\quad I=R^{>0},\quad
 H_G^*(X;\Z)=H^*(EG\times_GX;\Z),
\]
with the analogous definition for \(Y\). Assume both Borel cohomologies are finitely generated graded \(R\)-modules and that the canonical restriction maps give compatible base-change isomorphisms
\[
 H_G^*(X;\Z)/I H_G^*(X;\Z)\cong H^*(X;\Z),\qquad
 H_G^*(Y;\Z)/I H_G^*(Y;\Z)\cong H^*(Y;\Z).
\]
The aid then reports equivariant surjectivity \(H_G^*(X;\Z)\to H_G^*(Y;\Z)\) by graded Nakayama. These finiteness and base-change assumptions are essential qualifications, not consequences asserted for the original singular level. The integral coefficient convention states the integral version of the reported bridge; its application to the historical problem must be checked.

Finally, on a complex manifold with complex structure \(J\), let \(A,B\) be complete commuting real-holomorphic vector fields with flows \(\Phi_s,\Psi_t\). For the specified normalization \(z=e^{s+it}\), the action
\[
 z\cdot x=\Phi_s\Psi_t(x)
\]
is reported to descend from \(\R^2\) to a holomorphic \(\C^*\)-action exactly when
\[
 B=JA,\qquad \Psi_{2\pi}=\operatorname{id}.
\]
No effectiveness or minimal period of \(2\pi\) is required. Completeness and commutativity alone do not provide this complex action; the source does not verify these extra conditions for the moment-map flows at issue.

As dated context, the aid reports rigorous published results for linear torus actions and quiver-variety surjectivity by a different method, an unverified historical singular-level [CG83] application, and withdrawal of a 2014 linear-type preprint for a crucial proof error. It supplies no complete references for those later status claims.
\subsection{Short English statement}
These comments record gaps and tentative arguments surrounding hyperkähler Kirwan surjectivity. The supplied aid offers conditional bridges from split homology to cohomology, from ordinary to equivariant cohomology, and from two real flows to a holomorphic multiplicative-group action. The required hypotheses are not verified for the old singular-level argument, and this record is not an independent conjecture.
\subsection{Sources}
\begin{itemize}
\item Local UnsolvedMath registry, permanent ID 20001950, exact source code catalogue entry; canonical \texttt{aim-geometry-notes.json}, \texttt{notes[287]}, tag \texttt{conjecture}; AIM Workshop Problem Lists, Geometry. Created and updated 14 August 2026.
\item \emph{Moment maps and surjectivity in various geometries}, AIM workshop 9--13 August 2004, Comments 1.3--1.5; \url{https://aimath.org/WWN/momentmaps/momentmaps.pdf}. The supplied verification assigns the initial Comment 1.2 fragment to AIM-GEOMETRY-0287 and the trailing 3-Sasakian heading to the next section. That source-verification claim is supplied evidence; the PDF was not accessed here.
\item Kirwan's manuscript, Sections 3 and 5; Sjamaar [Sja98]; Carrell--Goresky [CG83]. Only these names, labels, and contextual descriptions are supplied, without complete bibliographic entries. S. Wu and J. Weitsman's historical discussion is explicitly tentative.
\item Supplied attempt 1: \texttt{context\_only}, problem status at run \texttt{not\_a\_problem}, rigor label \texttt{proved}, novelty confidence \texttt{low}; classification \texttt{OPEN-TRIAGE}, suggested difficulty L2, registry status \texttt{open}. Candidate: conditional bridge lemmas. Later published-status and preprint-withdrawal claims lack full references in this record and require independent verification.
\item Ulam AI / UnsolvedMath Contributors; CC BY 4.0. Machine-generated research aid, not peer reviewed; independently verify before citation.
\end{itemize}
\end{document}
